\documentclass{article}
\usepackage[UTF8]{ctex}
\usepackage{amsmath,amssymb,amsfonts}
\usepackage{enumitem}
\usepackage{geometry}
\geometry{a4paper, margin=1in}
% === 错题本排版优化（不参与题目标识） ===
\setlength{\parindent}{0pt}
\setlength{\parskip}{0.35em plus 0.10em minus 0.08em}
\setlength{\emergencystretch}{3em}
\setlist{itemsep=0.25em, topsep=0.40em, parsep=0pt}
\allowdisplaybreaks[2]
% === 错题本排版优化结束 ===
\begin{document}
\section*{西安交通大学2025年数学分析真题}
\section*{180．西安交通大学2025年数学分析试卷}
\begin{itemize}

\subsection{题目1}
  \item[1．] 设 \(x_{n}>0, \lim _{n \rightarrow \infty} x_{n}=1\) ，则 \(\lim _{n \rightarrow \infty} \sqrt[n]{x_{1} x_{2} \cdots x_{n}}=\) \(\_\_\_\_\)．

\subsection{题目2}
  \item[2．] 设 \(f(x)=x^{2} e^{x}\) ，则 \(f^{(10)}(0)=\) \(\_\_\_\_\)．
\end{itemize}

\subsection{题目3}
3． \(\lim _{x \rightarrow 0}\left(\frac{2^{x^{2}}+1}{2^{x}+1}\right)^{\frac{1}{x}}=\)

\subsection{题目4}
4． \(\int_{0}^{\frac{\pi}{2}} \ln \sin x \mathrm{~d} x=\) \(\_\_\_\_\)．

\subsection{题目5}
5．函数 \(f(x, y)=x^{3}-x^{2}-x y+y^{2}\) 的极值点为 \(\_\_\_\_\)．

\subsection{题目6}
6．已知

\[
f(x, y)= \begin{cases}x y \frac{x^{2}-y^{2}}{x^{2}+y^{2}}, & (x, y) \neq(0,0) \\ 0, & (x, y)=(0,0)\end{cases}
\]

则 \(\frac{\partial^{2} f}{\partial x \partial y}(0,0)+\frac{\partial^{2} f}{\partial y \partial x}(0,0)=\) \(\_\_\_\_\)．

\subsection{题目7}
7．Laplace 算子 \(\Delta f=\frac{\partial^{2} f}{\partial x^{2}}+\frac{\partial^{2} f}{\partial y^{2}}\) 在极坐标变换 \(\left\{\begin{array}{l}x=r \cos \theta \\ y=r \sin \theta\end{array} \quad(0 \leq \theta<2 \pi)\right.\) 下的表达式为 \(\_\_\_\_\)．\\
\(\_\_\_\_\)

\subsection{题目8}
8．\(\left(\partial \mathbb{Q}^{2}\right)^{\circ}=\) \(\_\_\_\_\)．

\subsection{题目9}
9． \(\iint_{x^{2}+y^{2} \leq 4}\left(x^{2} \sin y+x \ln \left(1+y^{2}\right)\right) \mathrm{d} x \mathrm{~d} y=\)

\subsection{题目10}
10．设 \(C\) 为心脏线 \(r=1+\cos \theta\) ，积分 \(\oint_{C}\left(x+y^{2}\right) \mathrm{d} x+\left(y+x^{2}\right) \mathrm{d} y=\) \(\_\_\_\_\)．

\subsection{题目11}
1．设 \(f(x)\) 在 \([0,+\infty)\) 上一致连续，证明：\(f(x)=O(\hat{x})(x \rightarrow+\infty)\) ．

\subsection{题目12}
2．设幂级数 \(\sum_{n=0}^{\infty} a_{n} x^{n}\) 在 \((-1,1)\) 上收敛，且 \(a_{n} \geq 0(\forall n \geq 0)\) 。又设

\[
\lim _{x \rightarrow 1^{-}} \sum_{n=0}^{\infty} a_{n} x^{n}=S \in \mathbb{R}
\]

证明：\(\sum_{n=0}^{\infty} a_{n}\) 收敛于 \(S\) 。

\subsection{题目13}
3．设 \(x_{n}=1+\frac{1}{\sqrt{2}}+\cdots+\frac{1}{\sqrt{n}}-2 \sqrt{n}(\forall n \geq 1)\) ，证明：\(\left\{x_{n}\right\}\) 收敛。

\subsection{题目14}
4．设 \(\Gamma(a)=\int_{0}^{+\infty} x^{a-1} e^{-x} \mathrm{~d} x, a>0\) ．证明：

\begin{itemize}
  \item[（1）] \(\Gamma(a+1)=a \Gamma(a), a>0\) ．
  \item[] \((2) \ln \Gamma(a)\) 是凸函数．
\end{itemize}

\subsection{题目15}
5．设 \(E\) 是 \(\mathbb{R}^{n}\) 中的开集，函数 \(f: E \rightarrow \mathbb{R}^{m}\) 与 \(g: E \rightarrow \mathbb{R}^{m}\) 均可微．设 \(h(\mathbf{x})=\langle f(\mathbf{x}), g(\mathbf{x})\rangle\) 为 \(f(\mathbf{x})\)与 \(g(\mathbf{x})\) 的内积，证明：\(h(\mathbf{x})\) 在 \(E\) 上可微，且对任意的 \(\mathbf{x} \in E\) ，有

\[
h^{\prime}(\mathbf{x})=f(\mathbf{x})^{T} g^{\prime}(\mathbf{x})+g(\mathbf{x})^{T} f^{\prime}(\mathbf{x}) .
\]

\subsection{题目16}
6．设 \(\Omega \subset \mathbb{R}^{2}\) 为非空的连通开集，\(f: \Omega \rightarrow \mathbb{R}\) 为连续函数．且

\[
f(\mathbf{x})=\frac{1}{\pi r^{2}} \int_{B_{\mathbf{r}}(\mathbf{w})} f, \sqrt[4]{B_{\mathbf{r}}(\mathbf{x})} \subset \Omega
\]

这里 \(B_{r}(x)\) 表示以 \(x\) 为中心，半径为 \(r\) 的开圆盘．证明：若 \(f\) 在 \(\Omega\) 上有最大值，则 \(f\) 为常值函数．
\end{document}
